\subsection*{Parameter Passing between Croqtile and C++ Function}

Passing arguments between Croqtile and C++ requires the \texttt{choreo.h} header. Combine a raw C++ pointer with its dimensions to construct Croqtile spanned data, as shown below.

\begin{lstlisting}[language=croqtile,caption={Parameter Passing between Croqtile and C Function}]
#include "choreo.h"
void bar(const float* data, unsigned size) {}
__co__ void foo(f32 mdspan<2> d) {
  parallel p by 6 {
    call bar(d, |d|);     // Call the C++ function
  }
}
void foobar(float* a) {
  foo(croq::make_spanview<2>(a, {1, 2}));
}
\end{lstlisting}

The \texttt{make\_spanview} utility template wraps the data without copying it. In the Croqtile function \texttt{foo}, the corresponding parameter \texttt{d} has type \texttt{f32 mdspan\textless{}2\textgreater{}}. An implicit conversion allows \texttt{d} to be passed as the first argument to the C++ function \texttt{bar}; \texttt{|d|} supplies its total element count.

Scalar values can also pass between C++ and Croqtile functions. An \texttt{ituple}, however, is used only inside a Croqtile function.
